\documentclass[11pt]{article}
\usepackage[a4paper,margin=27mm]{geometry}
\usepackage[T1]{fontenc}
\usepackage{lmodern,microtype}
\usepackage{amsmath,amssymb,amsthm,mathtools,mathrsfs}
\usepackage{booktabs,enumitem}
\usepackage[hidelinks]{hyperref}
\hypersetup{pdftitle={prb_34ecf5fa262f4a85a0fbe4af9837c46b: merged partial research contribution},pdfauthor={OpenAI GPT-6 assisted research draft},pdfsubject={Partial mathematical results; original problem remains unresolved}}
\usepackage{fancyhdr}
\pagestyle{fancy}
\fancyhf{}
\fancyhead[L]{\small prb_34ecf5fa262f4a85a0fbe4af9837c46b}
\fancyhead[R]{\small Nongeneric stability: partial results}
\fancyfoot[C]{\thepage}
\setlength{\headheight}{14pt}
\setlength{\parskip}{3pt}
\setlength{\emergencystretch}{2em}
\setlist{itemsep=2pt,topsep=5pt}
\newtheorem{theorem}{Theorem}
\newtheorem{proposition}[theorem]{Proposition}
\newtheorem{lemma}[theorem]{Lemma}
\newtheorem{corollary}[theorem]{Corollary}
\theoremstyle{definition}
\newtheorem{definition}[theorem]{Definition}
\newtheorem{remark}[theorem]{Remark}
\newcommand{\NE}{\operatorname{NE}}
\newcommand{\QP}{\operatorname{QP}}
\newcommand{\Stab}{\operatorname{Stab}_0}
\newcommand{\CC}{\mathscr C}
\newcommand{\A}{\mathcal A}
\newcommand{\Up}{\mathcal U_0}
\newcommand{\R}{\mathbb R}
\newcommand{\Q}{\mathbb Q}
\newcommand{\N}{\mathbb N}
\newcommand{\supp}{\operatorname{supp}}
\newcommand{\dist}{\operatorname{dist}}
\newcommand{\im}{\operatorname{im}}
\DeclareMathOperator{\relint}{relint}
\title{Axiomatic containment of Mertens-stable sets without genericity:\\
\large Merged partial contribution and proof audit}
\author{OpenAI GPT-6 assisted research draft}
\date{18 October 2026 \quad\textbullet\quad Merged contribution, version 3}
\begin{document}
\maketitle
\thispagestyle{plain}

\section*{Contribution metadata}
\noindent\textbf{Permanent problem identifier:} \texttt{prb\_34ecf5fa262f4a85a0fbe4af9837c46b}.\\
\textbf{Claim type:} \textbf{Partial result}, with full written proof arguments
and an exact global reduction; \textbf{not a complete solution}.\\
\textbf{Scope:} Largest globally axiom-compatible solution rule by compact
pruning; an exact reformulation of the missing step; a continuous-section
criterion (including fully mixed and strict equilibria); a nongeneric
zero-sum polymatrix family; duplicate-action and strong-embedding transfers;
two payoff-limit counterexamples; and a polynomial Hopf-map obstruction to
replacing rational-cohomological essentiality by a coincidence property.\\
\textbf{Assumptions and limitations:} Exact axioms $B,D^*,I^*$ and
Mertens $0$-stability in Govindan--Wilson (2026). No proof or globally
compatible counterexample for the unrestricted implication; no novelty,
publication, independent referee, or formal-verification claim.\\
\textbf{Source relation:} Consolidates the two user-supplied local drafts
for \texttt{prb\_34ecf5fa262f4a85a0fbe4af9837c46b}. The source author reports proof and boundary checks
against the pinned primary-source definitions; these reports are author
declarations, not independent verification.

\section{Write-up}
\paragraph{Executive statement.}
The generic-game implication is due to Govindan--Wilson. The precise
nongeneric implication is still unresolved. This manuscript gives several
sufficient subclasses, a global reduction, and explicit obstructions to
plausible but invalid extensions of the generic proof. The constructions
are deliberately not described as a universal solution.

\subsection{Exact target and external inputs}
A finite normal-form game is $G=(I,(A_i)_{i\in I},(u_i)_{i\in I})$,
with mixed-strategy space $\Sigma_G=\prod_i\Delta(A_i)$ and Nash set
$\NE(G)$. All games have at least one player and nonempty finite action sets.
We use canonically labelled finite player and action sets; this makes the
collection of games with real payoff arrays a set, not a proper class.
Relabellings and all permitted representations remain included in the axioms.

The source's definition of Mertens $0$-stability, including its perturbation
space, boundary conventions and connexity requirement, is normative
\cite[Definition 2.4]{gw}. Here $0$ denotes characteristic zero, not a local
fixed-point index. The problem is the implication
\begin{equation}\label{eq:target}
 \Phi\text{ satisfies }B,D^*,I^*
 \quad\Longrightarrow\quad
 \forall G\ \forall K\in\Phi(G)\ \exists S\in\Stab(G): S\subseteq K.
\end{equation}
The selected sets are nonempty, closed and connected subsets of $\NE(G)$.
No additional closedness of the \emph{family} $\Phi(G)$ is assumed.

We retain precisely the source's permitted comparisons. Axiom $B$ requires
intersection with profiles equivalent to quasi-perfect equilibria for every
perfect-recall extensive-form representation. Axiom $D^*$ compares the source's
strongly irrelevant dominated-strategy extensions, which have the same Nash
set under the natural inclusion. Axiom $I^*$ requires equality between a
solution family and the family of images of solutions under every strong
embedding. Arbitrary payoff perturbations and arbitrary dominated-strategy
extensions are \emph{not} added to these axioms
\cite[Axioms 3.1, 4.7, 5.4; Definitions 4.4, 5.2]{gw}.

\paragraph{External game-theoretic input.}
We use the established all-game properties of the Mertens stability
rule: nonemptiness and satisfaction of $B,D^*,I^*$, as stated by
Govindan--Wilson and proved in their axioms discussion and Appendices A--B.
The \emph{separate} generic necessity conclusion of their Theorem 6.1
is not assumed for nongeneric games. The stable-image and stable-lift
properties under strong embeddings are in Appendix B. These literature theorems are inputs;
we do not claim to reprove or independently referee their full proofs.

\paragraph{What is new relative to the preceding draft.}
The original local attempt addressed strict quasi-perfect equilibria; the
user-supplied manuscript adds explicit fully mixed and strict-equilibrium
sections and a particularly economical nongeneric payoff-limit example.
The compact pruning arguments address global axiom quantifiers, while
Subsection \ref{sec:polymatrix} supplies games with only non-strict equilibria.
The full target \eqref{eq:target} remains explicitly separate.

\subsection{Compact spaces of candidate solutions}\label{sec:compact}
Write
\[
 \CC(G)=\{K\subseteq\NE(G):K\ne\varnothing,\ K\text{ compact and connected}\},
\]
with the Hausdorff metric $d_H$. Compactness and closedness of individual sets
are equivalent here because $\NE(G)$ is compact.

\begin{lemma}\label{lem:hyperspace}
For every finite game, $\CC(G)$ is compact. If
$f:\NE(H)\to\NE(G)$ is continuous, then
$f_*:\CC(H)\to\CC(G)$, defined by $f_*(L)=f(L)$, is continuous.
\end{lemma}
\begin{proof}
The Nash set is closed in the compact product of simplexes: each best-response
inequality is closed. For completeness, the hyperspace of nonempty compact
subsets of any compact metric space $M$ is compact. Finite nets in $M$ yield
finite nets in its hyperspace by taking all nonempty subsets of a finite net.
It is complete: for a Hausdorff-Cauchy sequence $K_j$, the functions
$z\mapsto\dist(z,K_j)$ converge uniformly to a continuous function $v$.
The zero set $K=\{v=0\}$ is nonempty by compactness of $M$. Uniform convergence
implies $\sup_{z\in K}\dist(z,K_j)\to0$. If the other Hausdorff distance did
not tend to zero, points in $K_j$ bounded away from $K$ would have a limit in
$M$; uniform convergence would put that limit in $K$, a contradiction.
Thus $K_j\to K$.

Connected subsets form a closed subspace of this hyperspace. Indeed, if a
Hausdorff limit $K$ of connected compact sets were disconnected, a separation
of $K$ into two nonempty compact sets would give disjoint open neighborhoods
with positive separation. All sufficiently close $K_j$ would lie in their
union and meet both, contradicting connectedness. This proves compactness of
$\CC(G)$. Finally, $f$ is uniformly continuous on its compact domain. A modulus
of uniform continuity bounds $d_H(f(L),f(L'))$ in terms of $d_H(L,L')$.
Continuous images preserve nonemptiness, compactness and connectedness.
\end{proof}

For a fixed perfect-recall representation $\Gamma$ of $G$, let
$Q_\Gamma\subseteq\NE(G)$ denote exactly the mixed profiles equivalent to
quasi-perfect equilibria in the source's sense.

\begin{lemma}\label{lem:qpcompact}
$Q_\Gamma$ is compact. Consequently
$\{K\in\CC(G):K\cap Q_\Gamma\ne\varnothing\}$ is closed in $\CC(G)$.
\end{lemma}
\begin{proof}
First show that quasi-perfect behavioral profiles form a closed set for a
fixed finite $\Gamma$. Suppose $b^j\to b$ and every $b^j$ is quasi-perfect.
There are finitely many information-set actions. After discarding finitely
many terms, every action with positive probability in $b$ also has positive
probability in $b^j$. From a defining completely mixed witness sequence for
$b^j$, choose $c^j$ with $\|c^j-b^j\|<1/j$. Every action supported by $b^j$
is optimal against $c^j_{-i}$ at its information set, by
\cite[Definition 2.1]{gw}. Therefore every action supported by $b$ has the same
property. As $c^j\to b$, this is a quasi-perfect witness for $b$.
The behavioral strategy space is compact, so its quasi-perfect subset is
compact.

The mixed/behavioral equivalence relation used here is closed. In the source's
realization or enabling representation it is expressed by equality of finitely
many continuous realization coordinates: they are linear in mixed strategies
and polynomial in behavioral strategies. Thus the set of pairs $(x,b)$ with
$b$ quasi-perfect and $x$ equivalent to $b$ is compact. Its projection to
mixed-strategy space is $Q_\Gamma$ and is compact. The usual Nash inequalities
are preserved by this equivalence.

Finally, if $K_j\to K$ and $x_j\in K_j\cap Q_\Gamma$, take a convergent
subsequence in compact $Q_\Gamma$. Its limit lies in $K$ by Hausdorff
convergence. Hence $K\cap Q_\Gamma\ne\varnothing$.
\end{proof}

\subsection{A largest axiom-compatible rule by compact pruning}\label{sec:pruning}
\subsubsection{Comparison maps}
Introduce one arrow $e:H\to G$ for each of the following \emph{exactly
permitted} comparisons:

\begin{enumerate}[label=(\roman*)]
\item A strong embedding of $G$ in $H$, with the map
$f_e:\NE(H)\to\NE(G)$ supplied by Definition 5.2 of the source.
\item An extension $H$ of $G$ satisfying Definition 4.4 of the source, with
$f_e$ the natural deletion of the unused added coordinates on the Nash sets.
\end{enumerate}
In (ii), the two Nash sets are naturally homeomorphic. In (i), Nash profiles
indeed map to Nash profiles: the outsiders of a Nash profile lie in the
specified face; surjectivity of each insider's strategy map implements any
original deviation; and the source's preference preservation would turn a
profitable original deviation into a profitable insider deviation in $H$.
The other insiders' mapped strategies remain unchanged because their maps
depend only on their own strategies and the outsiders' strategies.

Both kinds of $f_e$ are continuous. Put $e_*(L)=f_e(L)$. For these arrows,
the two invariance axioms become the single family of equalities
\begin{equation}\label{eq:arrows}
 \Phi(G)=e_*(\Phi(H))
 :=\{e_*(L):L\in\Phi(H)\}\qquad(e:H\to G).
\end{equation}
For a dominated-strategy arrow this is just the source's equality under the
natural face identification, not an added projection axiom.

\subsubsection{The recursion}
Let $\mathfrak E(G)$ be the set of all permitted perfect-recall representations
of $G$. Define
\begin{equation}\label{eq:A0}
 \A_0(G)=\bigcap_{\Gamma\in\mathfrak E(G)}
 \{K\in\CC(G):K\cap Q_\Gamma\ne\varnothing\}.
\end{equation}
Thus $\A_0$ imposes exactly $B$ on individual candidate sets. For $r\ge0$, define
\begin{equation}\label{eq:recursion}
\begin{split}
 \A_{r+1}(G)={}&\A_r(G)\ \cap\!
 \bigcap_{e:G\to F}e_*^{-1}(\A_r(F))\ \cap\!
 \bigcap_{e:H\to G}e_*(\A_r(H)),\\
 \A_\infty(G)={}&\bigcap_{r=0}^{\infty}\A_r(G).
\end{split}
\end{equation}
The first intersection of comparisons enforces image compatibility; the
second enforces the existence of a compatible lift for \emph{each} incoming
comparison. Different arrows may require different lifts, exactly as in
\eqref{eq:arrows}.

\begin{theorem}[Largest compatible solution concept]\label{thm:maximal}
For every $G$, all families $\A_r(G)$ and $\A_\infty(G)$ are nonempty compact
subsets of $\CC(G)$. The rule $\A_\infty$ satisfies $B,D^*,I^*$ globally.
Moreover, every globally defined solution concept $\Phi$ satisfying these
axioms obeys
\begin{equation}\label{eq:largest}
 \Phi(G)\subseteq\A_\infty(G)\qquad\text{for every finite game }G.
\end{equation}
Thus $\A_\infty$ is the largest such rule, ordered by inclusion of solution
families game by game.
\end{theorem}
\begin{proof}
\emph{Compactness.} Lemma \ref{lem:qpcompact} makes $\A_0(G)$ closed in compact
$\CC(G)$. Suppose all $\A_r(F)$ are compact. Each inverse image in
\eqref{eq:recursion} is closed by Lemma \ref{lem:hyperspace}. Each image
$e_*(\A_r(H))$ is compact, hence closed in the metric space $\CC(G)$.
Arbitrary intersections of these closed sets are closed. This proves
compactness by induction. The families decrease with $r$.

\emph{Nonemptiness.} We prove the stronger assertion
\begin{equation}\label{eq:stableprotected}
 \Stab(G)\subseteq\A_r(G)\qquad\text{for every }G,r.
\end{equation}
At $r=0$, it follows from stability satisfying $B$. Suppose it holds at stage
$r$. For a stable set $S$ in $G$, every outgoing arrow has stable image in its
target, hence an image in that target's $\A_r$. Every incoming arrow admits a
stable lift of $S$, hence a lift in its source's $\A_r$. These statements are
exactly the established $D^*,I^*$ properties of $\Stab$.
Thus $S\in\A_{r+1}(G)$. Existence of stable sets proves nonemptiness, including
at the intersection $\A_\infty(G)$.

\emph{The axioms at the limit.} Axiom $B$ holds because
$\A_\infty(G)\subseteq\A_0(G)$. Fix an arrow $e:H\to G$. For
$L\in\A_\infty(H)$, membership in $\A_{r+1}(H)$ implies
$e_*(L)\in\A_r(G)$ for every $r$. Hence
\[
 e_*(\A_\infty(H))\subseteq\A_\infty(G).
\]
For the converse fix $K\in\A_\infty(G)$. For every $r$, the fiber
\[
 F_r=\{L\in\A_r(H):e_*(L)=K\}
\]
is nonempty: use $K\in\A_{r+1}(G)$ and the incoming-arrow condition.
It is compact by continuity of $e_*$, and $F_{r+1}\subseteq F_r$.
Compactness gives a point in $\bigcap_rF_r$. Such a point is a set
$L\in\A_\infty(H)$ with $e_*(L)=K$. This proves the reverse inclusion and
therefore \eqref{eq:arrows}. It is the step that justifies stopping at the
countable intersection; independently choosing a different lift at each
stage would not suffice without compactness.

\emph{Maximality.} Let $\Phi$ satisfy the three axioms. Axiom $B$ gives
$\Phi(G)\subseteq\A_0(G)$ for every $G$. Suppose inclusion holds at stage $r$.
For $K\in\Phi(G)$, its image under every outgoing arrow belongs to the
corresponding $\Phi(F)$ and hence $\A_r(F)$. Under every incoming arrow it has
a lift in $\Phi(H)$ and hence $\A_r(H)$. Thus $K\in\A_{r+1}(G)$.
Induction and intersection prove \eqref{eq:largest}. No compactness assumption
on the original family $\Phi(G)$ has been used.
\end{proof}

\begin{remark}
The existence of a largest rule also follows by taking the union of all
compatible rules. The additional content of Theorem \ref{thm:maximal} is the
explicit compact recursion, the compactness of its families, and the fact that
countably many stages suffice. The recursion is set-theoretic: its individual
stages quantify over all permitted finite games and comparisons.
\end{remark}

\subsection{An exact reformulation of the unsolved implication}\label{sec:equivalence}
Define the upward containment family
\[
 \Up(G)=\{K\in\CC(G):\exists S\in\Stab(G),\ S\subseteq K\}.
\]
This definition does not say that $K$ itself is stable.

\begin{corollary}\label{cor:equivalence}
The following are equivalent:
\begin{enumerate}[label=(\alph*)]
\item The unrestricted implication \eqref{eq:target} holds.
\item $\A_\infty(G)\subseteq\Up(G)$ for every finite game $G$.
\item For every $G$ and every $K\in\CC(G)\setminus\Up(G)$, there is a finite
integer $r\ge0$ such that $K\notin\A_r(G)$.
\end{enumerate}
If an explicit $G,K$ violating (b) is found, $\Phi=\A_\infty$ supplies the
required globally defined counterexample solution concept.
\end{corollary}
\begin{proof}
Apply (a) to $\Phi=\A_\infty$ to obtain (b). Conversely, maximality in
Theorem \ref{thm:maximal} makes (b) imply (a) for every compatible $\Phi$.
Finally, (b) and (c) are equivalent by the definition
$\A_\infty(G)=\bigcap_{r\ge0}\A_r(G)$.
\end{proof}

\paragraph{What this does not prove.}
The protected inclusion \eqref{eq:stableprotected} has the opposite direction
from the desired containment assertion. We have \emph{not} proved that every
survivor contains a stable set. Nor is finite-stage elimination an effective
finite certificate: one stage already quantifies over an unrestricted set of
representations and comparisons. There is no algorithm, uniform depth bound,
or test of all incoming embeddings in this contribution.

For games covered by the source's generic theorem, Corollary
\ref{cor:equivalence} does imply finite-stage elimination of each candidate
containing no stable set. Removing the genericity assumption from that
consequence is exactly the remaining issue, not an additional established
result.

\subsection{A constructive criterion for singleton stability}\label{sec:section}
Fix $G$, write $m_i=|A_i|$, and use the common-mass perturbation polytope
\[
 P_\varepsilon=\{\delta\tau:0\le\delta\le\varepsilon,\ \tau\in\Sigma_G\}.
\]
For $\eta\in P_1$, its common mass is
$\delta(\eta)=\sum_{a\in A_i}\eta_i(a)$, independent of $i$.
The perturbed feasible set is
$\Sigma_i(\eta)=\{x_i\in\Delta(A_i):x_i(a)\ge\eta_i(a)\}$.
Let $E_G$ be the graph of constrained Nash equilibria over $P_1$ and
$\pi:E_G\to P_1$ its projection.

The polytope $P_\varepsilon$ has dimension
\[
 d=1+\sum_i(m_i-1)
\]
in its linear span and is homeomorphic, with its relative boundary, to
$(D^d,S^{d-1})$. For $X\subseteq E_G$ put
$X_\varepsilon=X\cap\pi^{-1}(P_\varepsilon)$ and
$\partial X_\varepsilon=X\cap\pi^{-1}(\partial P_\varepsilon)$.
In particular $\partial X_1$ refers to the boundary of $P_1$, not an
arbitrarily chosen boundary of $X$.

\begin{proposition}[Continuous-section criterion]\label{prop:section}
Let $x\in\NE(G)$. Suppose that for some $0<\rho<1$ there is a continuous map
$s:P_\rho\to\Sigma_G$ with
\[
 s(0)=x,\qquad s(\eta)\in\NE(G(\eta))\quad(\eta\in P_\rho).
\]
Then $\{x\}\in\Stab(G)$.
\end{proposition}
\begin{proof}
Let $X=\{(\eta,s(\eta)):\eta\in P_\rho\}$. This is a compact subset of $E_G$
with $X_0=\{(0,x)\}$. For every $0<\varepsilon\le\rho$, projection is a
homeomorphism of pairs
\[
 (X_\varepsilon,\partial X_\varepsilon)
 \longrightarrow(P_\varepsilon,\partial P_\varepsilon).
\]
Consequently its induced map on degree-$d$ relative \v Cech cohomology with
coefficients in $\Q$ is an isomorphism between one-dimensional vector spaces,
so it is nonzero.

We verify the source's connexity condition separately. Any neighborhood $V$
of $(0,x)$ in $X$ contains the graph over $P_\kappa$ for some
$0<\kappa<\rho$: continuity at zero gives this after shrinking $\kappa$.
The graph over
\[
 B_\kappa=\{\eta:0<\delta(\eta)<\kappa,\ \eta_i(a)>0\text{ for all }i,a\}
\]
is connected and lies in $V\setminus\partial X_1$.
Connectedness follows, for example, by parameterizing $B_\kappa$ by
$(0,\kappa)\times\relint\Sigma_G$. The closure of this graph contains the
graph over $P_{\kappa/2}$, since strictly positive directions are dense.
Hence the connected component of $V\setminus\partial X_1$ containing this
graph has closure containing a neighborhood of $(0,x)$ in $X$.
The same conclusion holds whether closure is taken relative to $V$ or to
$X$, since that smaller graph lies in $V$. This is precisely connexity.
Together with essentiality it proves the assertion.
\end{proof}

\begin{corollary}[Equilibrium segments]\label{cor:segment}
Suppose $x,y\in\NE(G)$, $y$ is completely mixed, and the whole segment
$[x,y]$ lies in $\NE(G)$. Then $\{x\}$ is stable.
Consequently, if $\NE(G)$ is star-shaped about a completely mixed equilibrium,
every singleton equilibrium is stable and \eqref{eq:target} holds for $G$
for \emph{any} rule selecting nonempty subsets of $\NE(G)$.
\end{corollary}
\begin{proof}
Set $\mu=\min_{i,a}y_i(a)>0$ and choose $0<\rho<\min\{\mu,1\}$.
For $\eta\in P_\rho$, define
\[
 s(\eta)=\left(1-\frac{\delta(\eta)}{\mu}\right)x+
                \frac{\delta(\eta)}{\mu}y.
\]
It is an original Nash equilibrium by the segment assumption. Every coordinate
is at least $\delta(\eta)y_i(a)/\mu\ge\delta(\eta)\ge\eta_i(a)$, so it is
feasible for the perturbed game. Restricting deviations cannot invalidate an
original equilibrium that remains feasible. The continuous-section criterion
applies. Every nonempty selected set contains a point $x$, hence its stable
singleton.
\end{proof}

\begin{corollary}[Strict and completely mixed equilibria]\label{cor:strict}
Every strict equilibrium and every completely mixed equilibrium has a stable
singleton, under the stipulated strategy-perturbation definition.
\end{corollary}
\begin{proof}
For a completely mixed equilibrium take the constant section for sufficiently
small perturbations, or put $x=y$ in Corollary \ref{cor:segment}.
For a strict equilibrium $a^*$, assign exactly $\eta_i(a)$ to every
$a\ne a_i^*$ and all remaining probability to $a_i^*$. For sufficiently small
$\rho$, continuity and the finitely many strict payoff gaps ensure that
$a_i^*$ is the unique best reply at every such profile. Thus every player
allocates all discretionary probability optimally. This is a continuous
constrained-equilibrium section through $a^*$.
\end{proof}

\begin{remark}
These assertions concern \emph{strategy lower bounds in a fixed game}, not
arbitrary perturbations of its payoff array. Section \ref{sec:payofflimit}
shows why that distinction matters. The section criterion is sufficient,
not a claimed characterization of all stable singleton equilibria.
\end{remark}

\begin{corollary}[Backward induction and the earlier subclass]\label{cor:Bsection}
Suppose, for one permitted representation $\Gamma$ of $G$, every point of
$Q_\Gamma$ admits a section as in Proposition \ref{prop:section}. Then every
solution concept satisfying $B$ obeys stable-set containment at $G$.
In particular, if every quasi-perfect equilibrium in the simultaneous-move
representation is strict, every selected connected set is a stable singleton.
\end{corollary}
\begin{proof}
Axiom $B$ puts some $x\in Q_\Gamma$ in any selected $K$; the section gives
$\{x\}\in\Stab(G)$. In the strict case, strict payoff gaps persist in a
neighborhood of $x$. Every equilibrium in that neighborhood must put
probability one on the same uniquely optimal actions, so $x$ is isolated in
$\NE(G)$. Its singleton is relatively open and closed in the Nash set.
Connectedness of $K$ containing $x$ forces $K=\{x\}$.
\end{proof}

This retains the earlier draft's example:
\[
 \begin{array}{c|ccc}&L&M&R\\\hline
 U&(2,2)&(0,0)&(0,0)\\
 D&(0,0)&(0,0)&(0,0)
 \end{array}.
\]
Writing $p=\Pr(U)$ and $q=\Pr(L)$, the payoff advantages of $U$ over $D$
and of $L$ over either alternative are $2q$ and $2p$. If $p>0$, equilibrium
requires $q=1$ and then $p=1$. If $p=0$, action $D$ can be optimal only
when $q=0$. Hence the Nash set is $(U,L)$ together with the segment
$\{(D,tM+(1-t)R):0\le t\le1\}$. Against every completely mixed opponent,
$D$ is strictly inferior to $U$, so no point on the segment is quasi-perfect.
The strict point $(U,L)$ is quasi-perfect by continuity of its payoff gaps.
Corollary \ref{cor:Bsection} therefore selects only its singleton.

\subsection{A whole family with non-isolated quasi-perfect equilibria}\label{sec:polymatrix}
Let $n\ge3$ and let $A\in\R^{n\times n}$ satisfy
\[
 A^{\mathsf T}=-A,\qquad A\mathbf 1=0.
\]
Player $i$ chooses $a_i\in\{0,1\}$ and receives
\begin{equation}\label{eq:skewgame}
 u_i(a)=a_i\sum_{j=1}^n A_{ij}a_j.
\end{equation}
This is a finite polymatrix game: each payoff is a sum of pairwise terms.
It is zero-sum because $a^{\mathsf T}Aa=0$ for every pure profile.
Write $x_i=\Pr(a_i=1)$, so the mixed-strategy space is $[0,1]^n$.

\begin{theorem}\label{thm:skewgame}
For every game \eqref{eq:skewgame},
\begin{equation}\label{eq:kernelNE}
 \NE(G)=\ker A\cap[0,1]^n.
\end{equation}
Every equilibrium is quasi-perfect in the canonical simultaneous-move
representation and is non-strict. Every equilibrium singleton is Mertens
$0$-stable. Hence stable-set containment holds for every nonempty selected
subset, without needing $B,D^*,I^*$.
\end{theorem}
\begin{proof}
The payoff advantage of action $1$ over $0$ is $d_i=(Ax)_i$.
At a Nash profile, comparison with the feasible mixture $1/2$ gives
$d_i(x_i-1/2)\ge0$. Yet
\[
 \sum_i d_i(x_i-1/2)
 =x^{\mathsf T}Ax-\tfrac12\mathbf1^{\mathsf T}Ax=0,
\]
using skew-symmetry and $A\mathbf1=0$. Thus every term is zero.
If $d_i>0$, optimality requires $x_i=1$, making the term $d_i/2>0$.
If $d_i<0$, it requires $x_i=0$, making the term $-d_i/2>0$.
Both contradict the zero sum. Hence $Ax=0$. Conversely, $Ax=0$ makes both
actions best replies for every player, proving \eqref{eq:kernelNE} and
non-strictness.

The equilibrium set is convex and contains the completely mixed point
$c=\frac12\mathbf1$. For any $x\in\NE(G)$, the profiles
$x^k=(1-1/k)x+(1/k)c$ are completely mixed equilibria converging to $x$.
Against every $x^k_{-i}$ both actions are optimal. This proves quasi-perfection
using its defining completely mixed witnesses.

For clarity, singleton stability also has a direct affine certificate.
Given $0<\rho<1/2$ and $\eta\in P_\rho$, set
\begin{equation}\label{eq:affinesection}
 s_x(\eta)=(1-2\delta(\eta))x+\delta(\eta)\mathbf1.
\end{equation}
It belongs to $\ker A$ and every coordinate is in
$[\delta(\eta),1-\delta(\eta)]$. Both action probabilities therefore meet
all lower bounds, each of which is at most the common mass $\delta(\eta)$.
It is an original Nash equilibrium and remains a constrained one.
Proposition \ref{prop:section} proves stability of $\{x\}$.
\end{proof}

\subsubsection{An explicit three-player instance without duplicate actions}
Take
\[
 A=\begin{pmatrix}0&1&-1\\-1&0&1\\1&-1&0\end{pmatrix},
 \qquad
 u_i(a)=a_i(a_{i+1}-a_{i+2})
\]
with cyclic indices. Its equilibrium set is exactly
\[
 D=\{(t,t,t):0\le t\le1\}.
\]
No point of $D$ is isolated, and none is strict. Both actions of each player
are genuinely distinct: their payoff difference is $1$ at one opponents'
pure profile and $-1$ at another. Thus this continuum is not merely created
by duplicating payoff-equivalent actions.

For every $t\in[0,1]$, formula \eqref{eq:affinesection} becomes
\[
 s_t(\eta)=\bigl(t+\delta(\eta)(1-2t)\bigr)\mathbf1.
\]
It explicitly handles both boundary equilibria $t=0,1$ as well as all
interior equilibria. Block-diagonal matrices containing $k$ copies of this
$3\times3$ matrix yield $3k$-player examples with a $k$-dimensional equilibrium
cube and no strict equilibria. The same proof and section apply.

\subsection{Two useful transfer results}
\subsubsection{Lifting a section through exact duplicates}
\begin{proposition}\label{prop:duplicates}
Suppose $H$ is obtained by exact duplication of pure actions of $G$.
If $x\in\NE(G)$ has a continuous lower-bound equilibrium section, then every
mixed-profile lift $\widetilde x$ of $x$ under the duplication map has such a
section in $H$. In particular, every singleton $\{\widetilde x\}$ is stable.
\end{proposition}
\begin{proof}
Let $q_i:B_i\twoheadrightarrow A_i$ identify duplicates, so
$u_i^H(b)=u_i^G(q(b))$. Given lower bounds $\widetilde\eta$, aggregate them as
$\eta_i(a)=\sum_{q_i(b)=a}\widetilde\eta_i(b)$. Common perturbation mass is
preserved. Write $s(\eta)$ for a section through $x$ in $G$.
For each group $q_i^{-1}(a)$ choose weights $\lambda_i(b)\ge0$ summing to $1$.
When $x_i(a)>0$, choose $\lambda_i(b)=\widetilde x_i(b)/x_i(a)$; when
$x_i(a)=0$, all corresponding $\widetilde x_i(b)$ are zero and arbitrary
weights suffice. Define
\[
 \widetilde s_i(b)=\widetilde\eta_i(b)+\lambda_i(b)
                  \bigl(s_i(a;\eta)-\eta_i(a)\bigr),\qquad a=q_i(b).
\]
These coordinates meet their lower bounds, sum to one, and aggregate to
$s_i(\eta)$. They are continuous and equal $\widetilde x$ at zero.
Any feasible deviation in $H$ aggregates to a feasible deviation in $G$;
its payoff is exactly the corresponding original payoff. Since $s(\eta)$
is a constrained equilibrium, no such deviation is profitable.
The section criterion completes the proof, including zero-probability
boundary groups.
\end{proof}

\subsubsection{Containment descends through strong embeddings}
\begin{proposition}\label{prop:descent}
Let $\Phi$ satisfy $I^*$ and let $H$ strongly embed $G$ by $f$. If every
$L\in\Phi(H)$ contains a stable subset, then every $K\in\Phi(G)$ contains a
stable subset. For $D^*$ extensions, containment transfers in both directions
under the natural identification.
\end{proposition}
\begin{proof}
Given $K\in\Phi(G)$, choose $L\in\Phi(H)$ with $f(L)=K$ by $I^*$.
By hypothesis, some $S_H\in\Stab(H)$ is contained in $L$. The established
stable-image property gives $f(S_H)\in\Stab(G)$, and $f(S_H)\subseteq K$.
For a $D^*$ extension the Nash sets, selected families, and stable families
are identified, so the conclusion is immediate in either direction.
\end{proof}

The reverse strong-embedding inference is not supplied by this proof.
A stable $S\subseteq f(L)$ has \emph{some} stable lift under the literature's
stable-lift property; that property does not place this lift inside the
\emph{particular} selected set $L$. Confusing these two quantifiers would
silently assume the missing containment.

\subsection{Why a limit of generic-game conclusions is insufficient}\label{sec:payofflimit}
\begin{proposition}[One-sided tie payoff-limit obstruction]\label{prop:limit2}
There exist two-player games $H_\varepsilon\to H_0$ and stable singleton
sets $S_\varepsilon\to S_0$ in Hausdorff distance such that $S_0$ is
not stable. The limiting game has a nontrivial Nash-equilibrium continuum.
\end{proposition}
\begin{proof}
Let
\[
H_\varepsilon=\begin{array}{c|cc}
 &L&R\\\hline
 T&(1,1)&(1,0)\\
 B&(1+\varepsilon,1)&(0,0)
\end{array},\qquad \varepsilon\ge0.
\]
For every $\varepsilon>0$, $(B,L)$ is strict: $B$ strictly beats $T$
against $L$, and $L$ strictly beats $R$ against $B$. Hence
$S_\varepsilon=\{(B,L)\}$ is $0$-stable by the strict-equilibrium
section, and $S_\varepsilon\to S_0=\{(B,L)\}$.

In $H_0$, player 2 strictly prefers $L$ independently of the opponent;
if player 2 chooses $R$ with probability $r$, player 1's payoff
advantage of $T$ over $B$ is exactly $r$. Thus
\[
 \NE(H_0)=\{(pT+(1-p)B,L):0\le p\le1\}.
\]
For a perturbation vector $\eta$ with every lower bound strictly
positive and $0<\delta(\eta)<1$, the constrained game has a unique
Nash profile
\[
 x_1(T)=1-\eta_1(B),\qquad x_2(L)=1-\eta_2(R).
\]
As $\eta\to0$ through such interior perturbations, these profiles
converge to $(T,L)$, not to $(B,L)$. To verify failure of the actual
connexity requirement, suppose $S_0$ had a stability witness $X$ with
$X_0=\{(0,(B,L))\}$. Choose $V$ a small relative neighborhood of
$X_0$ in $X$ so that no profile in $V$ places more than $1/4$ on $T$
and $\delta(\eta)<1/4$. There are no graph points in $V$ over
$\operatorname{relint}P_1$, because every such equilibrium places
$1-\eta_1(B)>3/4$ on $T$. Hence
$V\setminus\partial X_1=\varnothing$, violating connexity.
Therefore $S_0$ is not stable. A continuous section through $(T,L)$
exists by the displayed formula, so this failure is specific to the
other endpoint and does not imply that $H_0$ lacks stable sets.
\end{proof}

The preceding example is incorporated from the user-supplied
\texttt{prb\_34ecf5fa262f4a85a0fbe4af9837c46b-attempt.tex}. It refutes only the proposed passage to
payoff limits; it does not give a counterexample to the globally
quantified axiomatic problem. In particular, backward induction in the
simultaneous-move representation excludes the endpoint $(B,L)$.

\begin{proposition}\label{prop:limit}
There are games $G_k\to G_0$ and singleton stable equilibria $K_k\to K_0$
in Hausdorff distance such that $K_0$ contains no stable subset.
\end{proposition}
\begin{proof}
For $\theta\ge0$ consider
\[
 G_\theta=\begin{array}{c|cc}
      &L&R\\ \hline
 U&(1,1)&(0,0)\\
 D&(0,0)&(\theta,\theta)
 \end{array}.
\]
For every $\theta>0$, $(D,R)$ is strict, so $K_\theta=\{(D,R)\}$ is stable
by Corollary \ref{cor:strict}. As $\theta\downarrow0$, the payoff arrays
converge and these singleton sets have the same limit $K_0$.

We show directly that $K_0$ is not stable in $G_0$, using the stipulated
connexity condition. Suppose it had a witness $X\subseteq E_{G_0}$ with
$X_0=\{(0,(D,R))\}$. Let $p=\Pr(U)$ and $q=\Pr(L)$, and consider the
relative neighborhood
\[
 V=X\cap\{\delta(\eta)<1/4,\ p<1/4,\ q<1/4\}
\]
of $X_0$. If all coordinates of $\eta$ are strictly positive, then
$q\ge\eta_2(L)>0$ makes $U$ strictly preferable to $D$, and
$p\ge\eta_1(U)>0$ makes $L$ strictly preferable to $R$.
Every constrained equilibrium must therefore satisfy
\[
 p=1-\eta_1(D)>3/4,\qquad q=1-\eta_2(R)>3/4
\]
when $\delta(\eta)<1/4$. No such profile belongs to $V$.
Consequently every parameter occurring in $V$ lies on the relative boundary
of $P_1$: it has a zero coordinate, including the zero-mass case.
Thus $V\setminus\partial X_1$ is empty, contradicting connexity.
Hence $K_0$ is not stable. Its only nonempty subset is itself.
\end{proof}

The positive-$\theta$ games are already nondegenerate bimatrix games.
If a particular stronger genericity convention is desired, take generic
payoff arrays arbitrarily close to each $G_{1/k}$ within the open region
where $(D,R)$ remains strict. They still converge to $G_0$ and give the same
counterexample. This invalidates the unqualified step ``approximate by
generic games and pass stability to the limit.'' It does not refute
\eqref{eq:target}: no globally compatible rule selecting $K_0$ is constructed.
Indeed the simultaneous-move backward-induction requirement already excludes
this particular $K_0$.

\subsection{A polynomial topological obstruction to a second shortcut}\label{sec:hopf}
The source identifies a gap between coincidence properties and rational
cohomological essentiality in the nongeneric case \cite[Section 10]{gw}.
The following classical Hopf-map construction makes the distinction explicit.
It is not asserted to be a Nash perturbation graph.

Let $D^m$ be the closed Euclidean unit ball. Regard $D^4$ as the unit ball
in $\mathbb C^2$ and define the real polynomial map
\begin{equation}\label{eq:hopf}
 F(z_1,z_2)=\bigl(2\operatorname{Re}(z_1\overline z_2),\,
 2\operatorname{Im}(z_1\overline z_2),\,|z_1|^2-|z_2|^2\bigr).
\end{equation}
Its boundary restriction is a Hopf map, up to the harmless choice of
orientation. The non-null-homotopy of this map is a standard input;
see \cite[Example 4.51, p.~380; Section 4.B, p.~427]{hatcher}.

\begin{proposition}\label{prop:hopf}
The map $F:D^4\to D^3$ has a coincidence with every continuous map
$g:D^4\to D^3$. Nevertheless,
\[
 F^*:\check H^3(D^3,S^2;\Q)\longrightarrow
           \check H^3(D^4,S^3;\Q)
\]
is zero. The restriction of $F$ to any compact subpair
$(Y,Y\cap S^3)$ also induces the zero map in degree three.
\end{proposition}
\begin{proof}
Direct calculation gives
$\|F(z)\|^2=(|z_1|^2+|z_2|^2)^2$; hence
$F^{-1}(S^2)=S^3$. Suppose a continuous $g:D^4\to D^3$ has no
coincidence with $F$. Then
\[
 R(z)=\frac{F(z)-g(z)}{\|F(z)-g(z)\|}
\]
is a continuous map $D^4\to S^2$. On the boundary $S^3$, consider
\[
 h_t(z)=\frac{F(z)-t\,g(z)}{\|F(z)-t\,g(z)\|},\qquad 0\le t\le1.
\]
For $t<1$, the denominator is nonzero because $\|F(z)\|=1$ and
$\|t g(z)\|\le t<1$; for $t=1$ nonvanishing follows from the
no-coincidence hypothesis. The map $h_t$ is a homotopy from
$F|_{S^3}$ to $R|_{S^3}$. But $R|_{S^3}$ extends over $D^4$ and is
therefore null-homotopic, which contradicts the non-null-homotopy of
the Hopf map. Thus a coincidence exists for every $g$.

For the ball pairs, relative \v Cech cohomology agrees with ordinary
relative cohomology, giving
\[
 \check H^3(D^3,S^2;\Q)\cong\Q,\qquad
 \check H^3(D^4,S^3;\Q)=0.
\]
Therefore $F^*=0$. For a compact $Y\subseteq D^4$, restriction of the
map of pairs to $(Y,Y\cap S^3)$ factors through $(D^4,S^3)$;
functoriality shows that the induced degree-three map remains zero.
\end{proof}

The graph
$\{(F(z),z):z\in D^4\}\subseteq D^3\times D^4$ is compact and semialgebraic,
and projection to its first coordinate is $F$ under a homeomorphism.
Thus merely insisting that a map be a polynomially described coordinate
projection does not cure the obstruction.
The example rules out a general topological inference from coincidence to
rational essentiality. To turn it into a counterexample to prb_34ecf5fa262f4a85a0fbe4af9837c46b, one would
still need to realize suitable structure in an actual Nash perturbation
graph and prove compatibility with \emph{all} comparisons in $B,D^*,I^*$.
Neither requirement is established here.

\subsection{Exact outstanding assertion and proof audit}\label{sec:gap}
The general theorem would follow from
\begin{equation}\label{eq:gap}
 K\in\A_\infty(G)\quad\Longrightarrow\quad
 \text{a stability witness has nonempty zero fiber contained in }\{0\}\times K.
\end{equation}
Such a witness must satisfy both the source's connexity condition and nonzero
rational relative cohomology. The compactness argument constructs compatible
\emph{selected sets}; it does not construct this perturbation witness.
The continuous-section argument supplies witnesses in the stated subclasses,
not for every survivor. No line of the manuscript proves \eqref{eq:gap}.

\paragraph{Quantifiers checked.}
The construction uses a global domain and both directions of every comparison.
Lifts at the limit come from nested compact fibers, not from an unlicensed
exchange of existence and limits. Boundary lower bounds, zero-probability
duplicate groups, and the pure endpoints of the equilibrium continuum are
included. The equilibrium examples do not use genericity or isolation.
Changing games in Proposition \ref{prop:limit} is expressly distinguished
from fixed-game hyperspace limits in Theorem \ref{thm:maximal}.

\paragraph{Scope of the counterexamples.}
Proposition \ref{prop:limit} refutes payoff-limit closure of singleton stability.
Proposition \ref{prop:hopf} refutes an unrestricted topological implication.
Neither is a counterexample to the posted game-theoretic statement.

\paragraph{Checks and reproducibility.}
The accompanying short Python program checks the displayed polynomial
identities symbolically and the explicit section inequalities using exact
rational test data. Those checks supplement the proofs; they do not establish
the universal axioms, the Hopf-map theorem, or the missing assertion
\eqref{eq:gap}. The manuscript is not a Lean or other proof-assistant
formalization, and no independent referee verification is claimed.

\paragraph{Novelty and provenance.}
The source's definitions, known stability properties and generic theorem
remain credited to their authors. Hyperspace compactness, section arguments,
and the Hopf obstruction use standard mathematics. The compact pruning
formulation and explicit examples were developed in this attempt; no claim
that they are absent from the literature is made. The prior manuscripts reported assistance from GPT-6 Astra Pro;
the present consolidation used ChatGPT GPT-6, without an available build
identifier or independently verified deployment settings. The user requested
and directed the work; no human verification is attributed to them.
This version supersedes the earlier local draft as a research write-up,
not as a change to the website's problem status.

\paragraph{Recommended submission metadata.}
\texttt{Problem: prb_34ecf5fa262f4a85a0fbe4af9837c46b}. Claim type: \texttt{partial result}.
Scope: largest-compatible-rule construction, exact elimination reformulation,
continuous-section criterion, a nongeneric polymatrix family, and two
obstructions to proposed proof routes. Status of prb_34ecf5fa262f4a85a0fbe4af9837c46b: \texttt{unresolved}.
No public issue, repository edit, or solved-status update is made by preparing
this manuscript.

\section*{Models used}
This consolidation used ChatGPT GPT-6 to compare the two supplied
LaTeX attempts, scrutinize the arguments, and prepare this merged draft.
Exact deployment snapshot and reasoning settings were not independently
available. The supplied manuscript declares ``OpenAI GPT 6 Astra Pro'';
that is preserved as the supplied manuscript's attribution, not asserted
as this run's verified model setting. No Lean or other formal proof
assistant was used.

\section*{Human contributors and roles}
The user supplied the WS-labelled attempt and the earlier local draft,
requested combination, and specified the target. No human mathematician's
proof endorsement or referee review is claimed. Govindan and Wilson
retain sole credit for the cited prior theory, and Hatcher for the cited
topological background. The added combination is an LLM-assisted research
draft with unverified novelty.

\begin{thebibliography}{9}
\raggedright
\bibitem{gw}
Srihari Govindan and Robert Wilson.
\emph{Axiomatic Equilibrium Selection in Generic Extensive-Form Games}.
arXiv:2504.16908v2, 26 August 2026.
The PDF title reads \emph{Axiomatic Equilibrium Selection: The Case of Generic
Extensive Form Games}.
Definitions 2.1, 2.4, 2.5, 4.4, 5.2; Axioms 3.1, 4.7, 5.4;
Theorem 6.1; Section 10; Appendices A and B.
\url{https://arxiv.org/html/2504.16908v2}.
\url{https://arxiv.org/pdf/2504.16908v2}.

\bibitem{hatcher}
Allen Hatcher.
\emph{Algebraic Topology}.
Cambridge University Press, 2002; author-hosted electronic edition.
Example 4.51, p.~380; Section 4.B, p.~427; the relative cohomology of balls.
\url{https://pi.math.cornell.edu/~hatcher/AT/AT.pdf}.

\bibitem{attachment}
User-supplied \emph{Local witnesses for nongeneric stability and two
obstructions to extending the generic theorem}, 18 October 2026,
file \texttt{prb\_34ecf5fa262f4a85a0fbe4af9837c46b-attempt.tex} (supplied working draft; not a
peer-reviewed source).

\bibitem{problem}
We-Solve.
\emph{Axiomatic containment of Mertens-stable sets without genericity}.
Permanent public problem identifier \texttt{prb\_34ecf5fa262f4a85a0fbe4af9837c46b}.
Consulted predecessor statement revision
\texttt{stm\_32cd4f81d2514415920d3ba7e50d23af}.
\href{https://we-solve.github.io/api/v1/problems/prb_34ecf5fa262f4a85a0fbe4af9837c46b.json}{Permanent problem lookup}.
\href{https://we-solve.github.io/problems/axiomatic-containment-of-mertens-stable-sets-without-genericity/}{Published problem record}.
Prompt blob SHA: \texttt{f34db89188d9a2b326623f8ccb664f0956da7f4b}.
\end{thebibliography}
\end{document}
